% =====================================================================
% VERSION ENSAE, REDUCTION SFE (phase 4, 1er octobre 2026) : notes de synthese.
% Consignes MS 2025-2026 du service des stages : deux notes autonomes, francaise puis
% anglaise, d'une a deux pages chacune, en fin de fichier. La version de quatre pages
% chacune, avec la figure LUCY et la figure N1, reste celle de la version Institut
% (chapitres/20_notes_synthese.tex) et l'historique Git de ce fichier.
% Lu par main_ensae.tex seul : ne jamais le faire lire par main_institut.tex.
%
% Chaque nombre est une valeur publiee dans le corps ; la note n'en calcule aucun. Elle cite
% ses scripts comme une section, pour rester sous harnais.
% =====================================================================

\providecommand*{\notesec}[1]{\par\smallskip\noindent{\bfseries #1}\par}

\chapter*{Note de synthèse}
\addcontentsline{toc}{chapter}{Note de synthèse}
\markboth{Note de synthèse}{}

% SOURCES-SCRIPTS: 28 30 43 48 60 63 67 68 74 91 92 95 110 08b 89 99 65 66 76 38 130 119 120 134 136

\notesec{Contexte et objectif}
Comment le besoin de capital varie-t-il entre états de conformité supposés ? Nous calculons
cet écart, puis examinons ce qui en subsiste quand la direction de la contagion et les coûts par
pilier ne sont pas observés. Les montants restent conditionnels à ces hypothèses.
Le règlement européen DORA distingue cinq domaines, appelés ici piliers : gouvernance (P1),
incidents (P2), tests (P3), prestataires tiers (P4) et partage d'informations (P5). La charge de
risque opérationnel de la Formule Standard ne dépend pas directement de leur conformité.

Le risque mesuré est le risque opérationnel d'origine informatique de l'entité. Les pénalités réglementaires et le risque de souscription des polices cyber
qu'un assureur vend ne sont pas modélisés. Nous répartissons l'écart entre les cinq piliers et entre
quatre paramètres du modèle que la non-conformité modifie, que nous appelons canaux : la fréquence
des incidents, leur détection, leur propagation d'un pilier à l'autre et la dépendance à un
prestataire commun. Le résultat est un besoin de capital au sens de l'ORSA (évaluation interne des risques), non un SCR réglementaire.

Les montants sont en millions d'euros et valent pour le périmètre financier étudié ($21{,}6$ incidents par an), non pour un assureur ; l'ordre de grandeur d'une entité de $20\,000$~M\$ d'actifs figure avec les résultats.
Le niveau change d'un facteur $22$ selon la seule source de sévérité. Sous la construction additive et les sensibilités étudiées, nous interprétons l'écart entre états et l'ordre des piliers ; le niveau reste illustratif.

\notesec{Démarche}
La base SAS OpRisk Global fournit $570$ pertes informatiques d'institutions financières de 2000 à 2026. Une loi de Pareto généralisée, ajustée au-delà de $20{,}03$~M\euro{} sur $86$ excès, décrit les
pertes élevées (indice de queue $0{,}5954$). Le nombre annuel d'incidents suit une loi binomiale négative. Un incident naît sur
un pilier, puis se transmet aux autres selon une matrice d'influence dirigée, posée à dire d'expert
puis bornée.  Le
capital est le quantile à $99{,}5\,\%$ de la charge annuelle, calculé dans chaque état.

La conformité modifie les paramètres du calcul. Les parts de vecteurs d'attaque Hackmageddon
pondèrent le multiplicateur de fréquence. La détection règle la probabilité qu'une perte
élémentaire dépasse le seuil ; elle s'applique à chaque pilier touché. La propagation règle
la poursuite de la marche entre piliers. Enfin, une amorce chez le prestataire P4 peut toucher
plusieurs autres piliers par un choc commun. Le tableau donne les valeurs retenues.

Ces différences entre états sont des hypothèses de scénario. Les tests de sensibilité du
chapitre~\ref{chap:multietats} examinent leurs effets ; les multiplicateurs de détection
ne sont pas balayés. Le calcul distingue l'effet d'un canal isolé et celui de sa correction
quand les autres restent dégradés.

\begin{center}\footnotesize
\renewcommand{\arraystretch}{1.1}
\begin{tabular}{@{}l l r r r r@{}}
\toprule
\textbf{Canal} & \textbf{Paramètre} & \textbf{Conforme} & \textbf{Non conforme} &
\textbf{Effet isolé} & \textbf{Effet de fermeture} \\
\midrule
Fréquence & incidents par an & $21{,}6$ & $53{,}6$ & $4\,328$ & $8\,775$ \\
Détection & dépassement par pilier & $0{,}128$ & $0{,}181$ & $1\,448$ & $4\,562$ \\
Propagation & gain de la cascade & $0{,}45$ & $0{,}90$ & $1\,633$ & $2\,402$ \\
Accumulation & prestataire commun & $0$ & $0{,}68$ & $1\,728$ & $3\,402$ \\
\midrule
\multicolumn{2}{@{}l}{\textbf{Besoin de capital} (M\euro)} & $6\,049$ & $20\,188$ & & \\
\bottomrule
\end{tabular}
\end{center}

\notesec{Résultats}
À l'échelle du périmètre financier étudié, le besoin de capital est multiplié par $3{,}34$ entre
l'entité conforme et l'entité non conforme sur les quatre canaux, soit un écart de $14\,139$~M\euro.
Les canaux pris un à un en expliquent $9\,138$~M\euro, et leurs interactions $5\,001$~M\euro, soit
$35\,\%$. L'effet isolé est le capital ajouté par un canal seul, depuis l'état conforme. L'effet de
fermeture est le capital retiré en corrigeant ce canal seul, depuis l'état non conforme sur les trois
autres : un plan de remédiation doit citer ce second effet.

\begin{center}
\includegraphics[width=0.78\textwidth]{N1_canaux_note.png}\par\vspace{1mm}
{\footnotesize Capital libéré en corrigeant un seul canal (M\euro, périmètre financier étudié) : en bleu
depuis une entité conforme ailleurs (effet isolé), en rouge depuis une entité défaillante partout
(effet de fermeture). Les effets de fermeture ne s'additionnent pas.\par}
\end{center}

La figure montre l'ordre des canaux. La fréquence domine dans les deux lectures. Corriger un canal
rapporte toujours davantage à une entité défaillante partout, car le canal ferme aussi les
interactions qu'il partageait avec les autres. Pour une entité, la priorité de remédiation dépend donc
de son état de départ, et non du seul effet isolé de chaque canal.

Cet écart dépend du coût de l'incident. À calage fixé, le calcul par Fourier donne
$13\,946$~M\euro{} sous la règle additive, à moins d'une erreur type ($367$~M\euro) de l'estimation
de Monte-Carlo, contre $5\,946$ avec une perte unique. Sous la règle additive, le bootstrap conjoint
de la sévérité donne une bande nominale à $95\,\%$ de $[4\,872\,;90\,361]$~M\euro{} pour l'écart,
mais de $[3{,}04\,;5{,}22]$ pour le facteur entre états (chapitre~\ref{ch:robustesse}). La queue
rend le niveau très incertain et le facteur beaucoup moins ; ces bandes restent conditionnelles à
la fréquence et aux autres paramètres posés.

Les données publiques ne disent pas quel pilier entraîne l'autre : nous calculons le capital aux
$1\,024$ sommets de l'ensemble des matrices d'influence admissibles, la structure symétrique étant
posée à dire d'expert. Dans le calage de référence (fréquence conforme, gain $0{,}90$), les valeurs aux sommets vont de $6\,858$ à
$8\,697$~M\euro{} ; recalculée sur quatre graines, la largeur de cette plage a un écart-type de
$189$~M\euro. La matrice se décompose en une partie symétrique, la force commune des liens,
posée ici, et une partie antisymétrique, leur sens, qu'aucune des sources
examinées ne fixe. Cette plage n'est ni un intervalle de confiance ni une borne démontrée : elle
mesure ce que coûte l'ignorance du sens, à force posée.  Pour une entité notionnelle de $20\,000$~M\$ d'actifs, le capital du secteur est
ramené à sa taille, la fréquence croissant avec les actifs selon un exposant estimé sur les grandes institutions : le besoin est de l'ordre de $24$~M\euro{} (de $15{,}8$ à $27{,}5$~M\euro{}
aux sommets admissibles), une extrapolation d'ordre de grandeur qui dépend de la structure d'expert et de la calibration.

\notesec{Limites et recommandations}
La fréquence de référence inclut douze pertes d'avant 2000, soit $2{,}1\,\%$ de plus que l'échantillon étudié.
Le chapitre~\ref{ch:robustesse} examine ce décalage.
La propagation et l'accumulation sont posées. Les correspondances testées conservent le signe
de l'écart sous la règle additive ; la perte par pilier reste non observée. Les bandes de calibration
laissent fixes la fréquence, les multiplicateurs et la cascade d'expert.

Le test du rapport de vraisemblance rejette la loi de Poisson (statistique de $130{,}5$) au profit de la
binomiale négative. Un backtest sur $2004$--$2025$ place ensuite $11$
années sur $12$ dans l'intervalle à $90\,\%$ de la loi binomiale négative, contre $9$ sur $12$ pour
Poisson, ce qui conforte le choix sans le démontrer sur douze années. Le backtest rejette le niveau de la charge annuelle ; une dérive de l'échelle de queue est explorée, sans confirmation statistique. Détecter un doublement du taux de dépassement avec $80\,\%$ de puissance demanderait $1\,811$ années à $99{,}5\,\%$.

Le modèle sert d'abord à comparer des scénarios et leurs effets de fermeture. Arrêter un capital
ou un investissement réel exige de mesurer les états, les coûts et l'exposition de l'entité.
Un registre daté, classé par domaine de contrôle et consignant tous les piliers touchés renseignerait
la direction. VERIS prévoit ces champs ; ils sont renseignés ensemble dans un incident sur $10\,591$.

% =====================================================================

\chapter*{Executive summary}
\addcontentsline{toc}{chapter}{Executive summary}
\markboth{Executive summary}{}

% SOURCES-SCRIPTS: 28 30 43 48 60 63 67 68 74 91 92 95 110 08b 89 99 65 66 76 38 130 119 120 134 136

\begin{otherlanguage*}{english}
\notesec{Context and objective}
How does capital need vary between assumed compliance states? We calculate this gap, then
examine what survives when contagion direction and pillar-level costs are unobserved. The amounts
remain conditional on these assumptions. The European DORA regulation distinguishes five areas,
called pillars here: governance (P1), incidents (P2), testing (P3), third-party providers (P4) and
information sharing (P5). The Standard Formula operational risk charge does not directly depend
on their compliance.

The risk measured is the entity's own IT-related operational risk, that is its own incidents.
Regulatory penalties and the underwriting risk of the cyber policies an insurer sells are not
modelled. We split the gap across the five pillars and across four model parameters that
non-compliance changes, which we call channels: incident frequency, detection, propagation from one
pillar to another, and dependence on a shared provider. The result is a capital need in the ORSA
sense, the internal assessment of risks and solvency. It is not a regulatory SCR.

The amounts are in millions of euros and are for the financial sample taken as a whole
($21.6$ incidents a year in the loss database), not for one insurer. The order of magnitude for an
entity with $20\,000$~M\$ of assets is given with the results. The level changes by a factor of $22$ with the severity source alone. Under the additive construction and the sensitivities examined, we interpret the gap between states and the ordering of pillars; the level remains illustrative.

\notesec{Approach}
SAS OpRisk Global provides $570$ IT losses at financial institutions from 2000 to 2026. A generalised Pareto distribution, fitted above $20.03$~M\euro{} on $86$ exceedances, describes
the high losses (tail index $0.5954$). The annual
incident count is negative binomial. An incident starts in one pillar, then spreads to the others
through a directed influence matrix, posited by expert judgement and then bounded.  Capital is the $99.5\,\%$ quantile
of the annual loss, computed in each state.

Compliance changes the calculation parameters. Hackmageddon attack-vector shares weight
the frequency multiplier. Detection sets the probability that an elementary loss exceeds
the threshold, applied to each affected pillar. Propagation sets the probability of continuing
the walk between pillars. Finally, an incident starting at provider pillar P4 can affect several
other pillars through a common shock. The table gives the values used.

These differences between states are scenario assumptions. The sensitivity tests in
Chapter~\ref{chap:multietats} examine their effects; detection multipliers are not swept.
The calculation distinguishes a standalone channel effect from the effect of correcting
that channel while the others remain degraded.

\begin{center}\footnotesize
\renewcommand{\arraystretch}{1.1}
\begin{tabular}{@{}l l r r r r@{}}
\toprule
\textbf{Channel} & \textbf{Parameter} & \textbf{Compliant} & \textbf{Non-compliant} &
\textbf{Standalone} & \textbf{Closure} \\
\midrule
Frequency & incidents per year & $21.6$ & $53.6$ & $4\,328$ & $8\,775$ \\
Detection & exceedance share & $0.128$ & $0.181$ & $1\,448$ & $4\,562$ \\
Propagation & cascade gain & $0.45$ & $0.90$ & $1\,633$ & $2\,402$ \\
Accumulation & shared provider & $0$ & $0.68$ & $1\,728$ & $3\,402$ \\
\midrule
\multicolumn{2}{@{}l}{\textbf{Capital need} (M\euro)} & $6\,049$ & $20\,188$ & & \\
\bottomrule
\end{tabular}
\end{center}

\notesec{Results}
At the scale of the financial sample, capital need is multiplied by $3.34$ between the
compliant entity and the entity that is non-compliant on all four channels, a gap of $14\,139$~M\euro.
The channels taken one by one explain $9\,138$~M\euro, and their interactions $5\,001$~M\euro, or
$35\,\%$. The standalone effect is the capital added by one channel alone, from the compliant state. The
closure effect is the capital removed by fixing that channel alone, from the state that is
non-compliant on the other three: a remediation plan should quote the second.

\begin{center}
\includegraphics[width=0.78\textwidth]{N1_canaux_note_en.png}\par\vspace{1mm}
{\footnotesize Capital released by fixing a single channel (M\euro, financial sample): in blue from an
entity compliant elsewhere (standalone effect), in red from an entity failing everywhere (closure
effect). Closure effects do not add up.\par}
\end{center}

The figure shows the ranking of channels. Frequency dominates under both readings. Fixing a channel
always releases more capital for an entity failing everywhere, because the channel also closes the
interactions it shared with the others. For an entity, the remediation priority therefore depends on
its starting state, not only on the standalone effect of each channel.

The gap depends on incident cost. At fixed calibration, Fourier computation gives
$13\,946$~M\euro{} under the additive rule, within one standard error ($367$~M\euro) of the
Monte Carlo estimate, against $5\,946$ with a single incident loss. Under
the additive rule, the joint bootstrap of severity gives a nominal $95\,\%$ band of
$[4\,872;90\,361]$~M\euro{} for the gap, but of $[3.04;5.22]$ for the factor between states
(Chapter~\ref{ch:robustesse}). The tail makes the level highly uncertain and the factor much less
so; both bands remain conditional on frequency and on the other posited parameters.

Public data do not say which pillar drives the other: we compute capital at the $1\,024$
vertices of the set of admissible influence matrices, the symmetric structure being set by expert
judgement. In the reference calibration (compliant-state frequency, gain $0.90$),
the vertex values range from $6\,858$ to $8\,697$~M\euro; recomputed over four seeds, the width
of this range has a standard deviation of $189$~M\euro.

The matrix splits into a symmetric part,
the common strength of the links, posited here, and an antisymmetric part, their
direction, which none of the sources examined pins down. This range is neither a confidence interval
nor a proven bound: it measures the cost of not knowing the direction, at posited strength.  For a notional entity with $20\,000$~M\$ of assets, sector capital is scaled down to its
size, frequency growing with assets at an exponent estimated on large institutions: the need is about $24$~M\euro{} ($15.8$ to $27.5$~M\euro{} at the admissible vertices),
an extrapolated order of magnitude that depends on the expert-set structure and on the calibration.

\notesec{Limitations and recommendations}
Reference frequency includes twelve losses from before 2000, $2.1\,\%$ more than
the calibration sample; Chapter~\ref{ch:robustesse} explains this convention.
Propagation and accumulation are posited. The tested mappings preserve the sign of the gap under
the additive rule; pillar-level loss remains unobserved. Calibration bands hold frequency,
multipliers and the expert-set cascade fixed.

A likelihood ratio test rejects the Poisson distribution (statistic of $130.5$) in favour of the
negative binomial. A backtest over $2004$--$2025$ then places $11$ years out of
$12$ inside the $90\,\%$ interval of the negative binomial, against $9$ out of $12$ for the Poisson.
This supports the choice without proving it over twelve years. The backtest rejects the annual aggregate level; a drift in tail scale is explored without statistical confirmation. Detecting a doubling of the exceedance rate with $80\,\%$ power would require $1\,811$ years at $99.5\,\%$.

The model primarily supports comparisons of scenarios and closure effects. Setting an actual
capital amount or investment requires measuring the entity's states, costs and exposure.
A dated register, classified by control domain and recording every pillar hit, would inform
direction. VERIS provides these fields; they are filled in together in only one incident out of
$10\,591$.
\end{otherlanguage*}
